\documentclass[12 pt]{article}
\usepackage[a4paper,
 mag=1000, includefoot,
 left=2cm, right=1.5cm, top=2cm, bottom=1cm, headsep=1cm, footskip=0.5cm]{geometry}
\usepackage[T2A]{fontenc}
\usepackage[utf8]{inputenc}
\usepackage[english,russian]{babel}
%\usepackage{indentfirst}
\usepackage{amssymb}
\usepackage{amsmath}
\usepackage{amsthm}
%\usepackage{array}
%\usepackage{graphicx}
%\usepackage{hyperref}
%\usepackage{url}
%\usepackage{caption}

\def\E{\mathbb{E}}
\def\P{\mathbb{P}}
\def\D{\mathbb{D}}

\begin{document}

\section*{Одномерный случай}
\subsection*{Распределения}

Нормальное распределение $N(a,\sigma^2)$ со средним $a$ и дисперсией $\sigma^2$. Свойства не перечисляю, надо знать. Также надо знать ЦПТ --- центральную предельную теорему.

Распределение хи-квадрат $\chi^2(m)$ с $m$ степенями свободы:\\
Пусть $\xi_i\sim N(0,1)$, независимы. Тогда $\eta_m = \sum_{i=1}^m \xi_i^2$
имеет распределение $\chi^2(m)$. $\E \eta_m = m$, $\D \eta_m = 2m$.\\
Распределение $(\eta_m -m )/\sqrt{2m}$ сходится к $N(0,1)$ при $m\rightarrow \infty$.

Распределение Стьюдента $t(m)$ с $m$ степенями свободы:\\
Пусть $\xi\sim N(0,1)$, $\eta_m\sim \chi^2(m)$, независимы. Тогда $\zeta_m = \xi/\sqrt{\eta_m/m}$
имеет распределение $t(m)$. $\E \zeta_m = 0$ ($m>1$), $\D \zeta_m = m/(m-2)$ ($m>2$).\\
Распределение $\zeta_m$ сходится к $N(0,1)$ при $m\rightarrow \infty$.

Распределение Фишера-Снедекора $F(m_1,m_2)$ с $m_1$ и $m_2$ степенями свободы:\\
Пусть $\eta_1\sim \chi^2(m_1)$, $\eta_2\sim \chi^2(m_2)$, независимы. Тогда $\phi_{m_1,m_2} = \frac{\eta_1/m_1}{\eta_2/m_2}$
имеет распределение $F(m_1,m_2)$.

Соотношение между распределениями:
$\xi^2 \sim \chi^2(1)$, $\zeta_m^2 \sim F(1,m)$, распределение $m_1\phi_{m_1,m_2}$ сходится к $\chi^2(m_1)$ при $m_2\rightarrow \infty$.

\subsection*{Оценки}
Пусть $\xi$ --- некоторая случайная величина с математическим ожиданием $a$ и дисперсией $\sigma^2$,
$(x_1,\ldots,x_n)$ --- выборка объема $n$.

$\bar{x} =\sum_{i=1}^n x_i/n$ --- несмещенная, состоятельная оценка $a$, $\E\bar{x} = a$, $\D\bar{x} = \sigma^2/n$.
Если $\xi\sim N(a,\sigma^2)$, то $\bar{x}\sim N(a, \sigma^2/n)$. В произвольном случае ($\D \xi < \infty$),
распределение $\sqrt{n}(\bar{x}-a) / \sigma$ стремится к $N(0, 1)$ (асимптотическая нормальность выборочного среднего).

$s^2=\sum_{i=1}^n (x_i-\bar{x})^2/n$ --- асимптотически несмещенная, состоятельная ($\D \xi < \infty$) оценка $\sigma^2$.
$\widetilde{s}^2=\sum_{i=1}^{n} (x_i-\bar{x})^2/(n-1)$ называется исправленной выборочной дисперсией,
$\sum_{i=1}^n (x_i-\bar{x})^2 = n s^2 = (n-1)\widetilde{s}^2$.
Если $\xi\sim N(a,\sigma^2)$, то $n s^2 / \sigma^2 = (n-1) \widetilde{s}^2 / \sigma^2 \sim \chi^2(n-1)$. \\
Иногда пишут: $\widetilde{s}^2 = \mbox{SS}/\mbox{d.f.}$, имея в виду, что SS --- sum of squares (сумма квадратов),
а d.f. --- число степеней свободы (degree of freedom) соответствующего распределения $\chi^2$.

\subsection*{Расстояние в сигмах}
В модели нормального распределения, расстояние от математического ожидания $a$ измеряется в сигмах $\sigma$.
`Далеко' --- то, что случается с маленькой вероятностью:\\
$\P(|\xi-a|>k\sigma) = 2(1-\Phi(k))$ = $\alpha$, где $\Phi$ --- функция распределения стандартного нормального распределения
$N(0,1)$.

Для $\alpha = 0.05$, $k=1.96\approx 2$, т.е. получаем правило двух сигм: вероятность попасть дальше, чем две сигмы, примерно равна 0.05 (мала).
Аналогично, правило трех сигм и пр.

\subsection*{Гипотезы и критерии}

\subsubsection*{Гипотеза $H_0: \E\xi = a_0$}

Статистика критерия (t-test): $t=\sqrt{n-1}(\bar{x}-a_0) / s = \sqrt{n}(\bar{x}-a_0) / \widetilde{s} $.
Если $\xi\sim N(a,\sigma^2)$, то $t$ имеет распределение Стьюдента $t(n-1)$. Иначе асимптотическое распределение
при $n\rightarrow \infty$ --- нормальное распределение $N(0,1)$. Также стандартное нормальное распределение получается
в нормальной модели с известной дисперсией и $t= \sqrt{n}(\bar{x}-a_0) / \sigma$.

Заметим, что квадрат статистики критерия распределен как $\chi^2(n-1)$ и может быть записан в форме\\
$t^2 = (\bar{x}-a_0)\left(\widetilde{s}^2/n\right)^{-1} (\bar{x}-a_0) = r^2(\bar{x},a_0)$.

Доверительный интервал в нормальной модели для $\E\xi$ имеет вид $\P(\E\xi \in (\bar{x}\pm c_\gamma \widetilde{s}/\sqrt{n}) = \gamma$,
где $t(n-1)(c_\gamma)=\gamma$. Также, доверительный интервал может быть записан в форме окрестности вокруг $\bar{x}$:
$\P(r^2(\bar{x},\E\xi) < c_\gamma^2)=\gamma$.

\subsubsection*{Гипотеза $H_0: \E\xi_1 = \E\xi_2$, признаки зависимые (измерены на одних и тех же индивидах)}
Данные имеют вид $(x_1,y_1)^\textrm{T}, \ldots, (x_n,y_n)^\textrm{T})$.
Эта гипотеза равносильна гипотезе $H_0:\E(\xi_1 -\xi_2) = 0$, т.е. сводится к предыдущему варианту для выборки
$z_i = x_i - y_i$.

\subsubsection*{Гипотеза $H_0: \E\xi_1 = \E\xi_2$, признаки независимые (один признак, измеренный на разных индивидах)}

Данные имеют вид $(x_1, \ldots, x_{n_1})$ и $(y_1, \ldots, y_{n_2})$.
Соответственно, есть $\bar{x}$, $\bar{y}$, $s_1^2$ и $s_2^2$, а также их исправленные варианты.

Для случая, когда предполагается, что $\D\xi_1 = \D\xi_2 = \sigma^2$,
введем в качестве оценки $\sigma^2$ понятие pooled variance (объединённой дисперсии, pool --- бассейн),
когда из каждой выборки вычитается свое среднее, а потом они объединятся (сливаются).
Введем ее исправленный вариант:\\
$\widetilde{s}^2 = \big((n_1-1) \widetilde{s}_1^2 + (n_2-1) \widetilde{s}_2^2\big)/ (n_1+n_2-2)
= \big(\sum_{i=1}^{n_1} (x_i-\bar{x})^2 + \sum_{i=1}^{n_2} (y_i-\bar{y})^2\big)/(n_1+n_2-2)$.

Если $\xi_i \sim N(a_i,\sigma^2)$, то $(n_1+n_2-2)\widetilde{s}^2/\sigma^2 \sim t(n_1+n_2-2)$.

Статистика критерия имеет вид\\
$ t = (\bar{x} - \bar{y})/\big(\widetilde{s}\sqrt{1/n_1 + 1/n_2}\big)$,
$ t^2 = (\bar{x} - \bar{y})\big(\widetilde{s}^2(1/n_1 + 1/n_2)\big)^{-1} (\bar{x} - \bar{y})$.

Если верна модель $\xi_i \sim N(a_i,\sigma^2)$, то $t$ имеет распределение Стьюдента $t(n_1+n_2 - 2)$,
а $t^2\sim F(1,n_1+n_2-2)$.
Иначе точное распределение неизвестно, но асимптотическое распределение при $n_1,n_2 \rightarrow \infty$
является стандартным нормальным.

\subsubsection*{Гипотеза $H_0: \E\xi_1 = \ldots = \E\xi_k$ для независимых выборок в нормальной модели}
Если модель имеет вид $\xi_i \sim N(a_i,\sigma^2)$ (т.е. еще и дисперсии одинаковые), то проверка этой
гипотезы называется дисперсионным анализом. Точнее --- одномерным однофакторным дисперсионным анализом
(1-way ANOVA, ANalysis Of VAriance).  Данные обычно обозначаются как
$y_{ij}$, где $i=1,\ldots,k$ --- номер группы индивидов, на которых измеряется  $\xi_i$,
$j=1,\ldots,n_i$ --- размеры групп, $n=\sum_{i=1}^k n_i$.  Статистика критерия строится на основе разложения дисперсии
(разложения суммы квадратов на выборочном языке) SStotal = SSbetween + SSwithin, где
SStotal $= \sum_{ij} (y_{ij} - \bar{y})^2\sim \chi^2(n-1)$, SSbetween $= \sum_{i} n_i (\bar{y}_i - \bar{y})^2 \sim \chi^2(k-1)$,
SSAwithin $= \sum_{ij} (y_{ij} - \bar{y}_i)^2 \sim \chi^2(n-k)$.
Статистика критерия имеет вид\\ $t=\Big(\mbox{SSbetween}/(k-1)\Big) / \Big(\mbox{SSwithin}/(n - k)\Big)\sim F(k-1, n-k)$.

\subsubsection*{Гипотеза $H_0:\D\xi_1 = \ldots = \D\xi_k$ для независимых выборок в нормальной модели}

Стандартный критерий Фишера для $k=2$ в нормальной модели: $t=\widetilde{s_1}^2/\widetilde{s_2}^2\sim F(n_1-1,n_2-2)$.
На гипотезу $H_0:\D\xi_1 = \ldots = \D\xi_k$ он обобщается как $t=\widetilde{s}_{\mbox{min}}^2/\widetilde{s}_{\mbox{max}}^2$
(и критическая область с одной стороны; распределение другое).
Есть еще критерий Левена, который тоже естественным образом обобщается на случай $k$ дисперсий.

Можно еще рассматривать критерий Бартлетта:\\
$t = A(n_1,n_2,k) \Big((\sum_{i=1}^k n_i -k)\ln \widetilde{s}^2 - \sum_{i=1}^k (n_i-1) \ln \widetilde{s_i}^2\Big)$
с асимптотическим распределением $\chi^2(k-1)$ при $\min n_i \rightarrow \infty$.

\newpage
\subsection*{Задания}
\begin{enumerate}
\item
Выборочное среднее равно 0.8, исправленная выборочная дисперсия равна 4, объем выборки равен 25. Проверьте гипотезу, что
мат.ожидание равно 2. Нужно найти p-level (вероятностный уровень) и сформулировать, при каких уровнях значимости
гипотеза отвергается, при каких не отвергается.
\item
Пусть $\eta\sim \chi^2(72)$. Можно ли приближенно считать расстояние от мат. ожидания в сигмах, почему?
Если да, то примерно в скольких сигмах находится 108 от 72?
\item
Проверить гипотезу, что скорость бега, в среднем, не зависит от настроения утром, если
в 10 дней с хорошем настроением скорость при утренней пробежке была, в среднем, 12 км/ч со стандартным отклонением 0.2,
тогда как за 20 дней с плохим настроением скорость бега была, в среднем, 13 км/ч со стандартным отклонением 0.5.
Проверьте гипотезу, условно предположив, что модель нормальная и дисперсии одинаковые.
\item
После года занятия в тренажерном зале у спортсмена возникла гипотеза, что за короткую тренировку давление не повышается.
В результате измерения давления в течение 25 дней до и после тренировки было получено, что давление, в среднем, вырастало на 1.
При этом стандартное отклонение до тренировки было равно 4, а после тренировки - 6. Корреляция между измерениями давления
оказалась равна 0.8.
\item
В трех группах с одинаковым числом студентов, равным 20,
по оценкам на экзамене были сосчитаны SStotal = 25, SSbetween=4. Проверьте гипотезу, что уровень студентов
во всех трех группах одинаковый.
\end{enumerate}
\end{document}
